\paragraph{Physical preparation and odd source response.}
A finite algebraic controller starts from physical vacua, constructs the
actual local ground doublets by polynomial projection, and prepares the
native cube trial with retained fermionic spin references. A Householder
unitary and a two-copy energy-agreement filter use the actual
$H_0+2^{-34}T$. The Cayley filter needs no exact interacting ground vector
or ground-energy input. With $1197721483502453176152084669$ terms,
preparation succeeds with probability at least $0.999999942661984$;
the conditional source trace distance from its true ground is at most
$2.355922487\,10^{-7}$. Every failure branch remains physical. Exact
algebraic synthesis, controlled matrix inversion and the finite schedule
are declared services; the construction supplies no efficient native circuit
(\path{model/coherent_defect_wave/chiral_joint/interacting_ground_preparation}, \cite{ophrepo}).

The separate finite native service compiler executes a two-state
singlet code of the actual $320$-state lepton/scalar source. Its exact
$HV=Vh$ has $h=\left(\begin{smallmatrix}0&-2\\-2&1\end{smallmatrix}\right)$,
so one simultaneous source pulse gives the full continuous code evolution;
six pointer pulses supply the complete native C/E readout. A compensated
$U(1)$ code also checks $491520$ CAR/Gauss columns and all $32768$
amplitudes. Rational control-error bounds retain export failures and
reset timeouts. Coherent load/export and signed analog port control are
explicit services; the passive classical repair run does not implement
them. The finite general compiler includes padding and an explicit
product-formula cost, without executing the full non-Abelian cube
preparation (\path{model/native_quantum_service_bridge}, \cite{ophrepo}).

On the thirty-mode reference vacuum, the fermion transfer is
\[
 B_r=(2d_r)^{-1/2}\sum_{s,a\in r}r^\dagger_{sa}c_{0sa},
 \qquad B_r^\dagger B_r=N_r/(2d_r)\le I.
\]
Its full dilation preserves
joint Gauss charges and total fermion parity. The reference energy is
$4N_{\rm ref}$; a bound over every reference occupation and source link
sector proves that the neutral ground with blank reference is the unique
physical ground of this enlarged Hamiltonian. Complete positive local
resolvents have nine energies for $e^c$ removal and seventeen for $L$
removal. Residual bounds transfer their weights to bins of the actual
full interacting Hamiltonian without assuming that unseen eigenstates
are absent from those bins:
\[
\begin{array}{c|c|c}
\text{native removal}&\text{source energy bin}&\text{selected probability}\\\hline
 e^c &[0.6906800974,0.6926332644]&[0.4197041521,0.4201989707]\\
 L &[0.0959119144,0.0978650814]&[0.3210643365,0.3214971263].
\end{array}
\]
The complete outside probabilities are enclosed respectively by
$[0.0799404942,0.0801565287]$ and $[0.1785579463,0.1788807366]$;
pump failure is about one half. The exact probabilities sum to one.
These ground-reference intervals acquire the stated preparation error
under the constructive protocol. Total source/reference excitation energy
adds the declared reference cost four. A localized charged response does
not identify an isolated propagating pole, a rest mass or an empirical
lepton; the static compensator is retained in that distinction
(\path{model/coherent_defect_wave/chiral_joint/interacting_charged_response}, \cite{ophrepo}).

\paragraph{Charge transport and its energy cost.}
The same prepared source and elementary removal have six orthogonal
neighbor-arrival routes. At source time $2^{-24}$ their unconditional
union probabilities exceed $3.8580\,10^{-37}$ for $e^c$ and
$6.2900\,10^{-37}$ for $L$. The amplitude estimate includes imperfect
preparation before squaring and retains every complementary outcome.
It does not identify those arrivals with the selected low energy bins.

For that distinction, let $\Pi_y$ require gauge-singlet matter at every
cell except $y$, retaining all scalar, fermion, link and reference states.
Gauss' law forces a nonblank link path from the origin reference to $y$.
With $\ell$ its graph distance, the complete local sector bounds imply
\[
 \Pi_y H_0\Pi_y\ge(8E_0+\ell+\delta_r)\Pi_y,
 \qquad \Pi_y T\Pi_y=0,
\]
where $\delta_{e^c}>0.10058529$ and $\delta_L>0.09688849$.
If $P_I$ is a selected full-Hamiltonian energy bin, with upper energy
$b<a_y=8E_0+\ell+\delta_r+4$, the separated-spectrum Sylvester identity
gives
\[
 \|\Pi_yP_I\|\le\frac{360\,2^{-34}}{a_y-b}.
\]
Thus at every source time the conditional probability of any of the
seven remote single-charge locations is below $8.656\,10^{-15}$ for
$e^c$ and $1.699\,10^{-15}$ for $L$. The eighth location and the
multi-charge cloud complete the location measurement; preparation,
pump and energy-selection failures give thirteen sequential records.
The cloud can contain screening. An explicit positive-norm state has a
locally compensated reference and a remote charge pair without a flux
path back to the origin. The bound therefore concerns the unscreened
single-charge domain, rather than unrestricted confinement
(\path{model/coherent_defect_wave/chiral_joint/interacting_charged_transport},
\path{model/coherent_defect_wave/chiral_joint/interacting_charge_localization},
\cite{ophrepo}).

The charged-particle criterion also depends on the intended field phase.
Under the relativistic Gauss-law and asymptotic-field hypotheses of
\href{https://doi.org/10.1016/0370-2693(86)91110-X}{Buchholz's infraparticle theorem},
a charged Coulomb sector has no sharp mass eigenstate. Its physical
mass target is a dressed or inclusive energy--momentum threshold with
specified detector resolution. The finite matrix bins above establish
neither that Coulomb phase nor a relativistic threshold. Demanding an
isolated charged pole in every continuum phase would impose an
inappropriate acceptance criterion.

\paragraph{Neutral propagation and a controlled field limit.}
The actual eight-cell ground $g$ and isolated triplet $e_\alpha$ supply
dressed native-current matrix units
$A_\alpha=P_gO_\alpha P_t/(60\sqrt r)=|g\rangle\langle e_\alpha|$,
where $r>0$ is the computed current residue and
$\Delta\in[8.20758,8.21035]10^{-22}$ is the source gap.
A declared number-conserving exchange between complete blocks gives
an isolated moving band
$E(k)=\Delta-2J\sum_i\cos k_i$ and a uniform positive gap.
Its controlled spacing limit is Schr\"odinger evolution.

A different, explicitly paired architecture places $N$ such blocks at
each cubic site. Write $A^{(N)}_{x\alpha}=N^{-1/2}\sum_j A_{xj\alpha}$
and $X=A^{(N)}+(A^{(N)})^\dagger$. On the full physical carrier declare
\[
 H_{a,N}=\frac1a\sum_{x,j}(h_{xj}-E_g)
       -t_a\sum_{\langle xy\rangle,\alpha}X_{x\alpha}X_{y\alpha},
 \qquad t_a=\frac{\Delta(1-a^2)}{12a},\quad 0<a\le1.
\]
Its clock is $e^{-itH_{a,N}}$. On the exact permutation-symmetric
sector, $(A^{(N)}_{x\alpha})^\dagger
=b^\dagger_{x\alpha}\sqrt{1-n_x/N}$, including the capacity boundary.
Every higher native sector is retained. A variational argument places a
full-source ground in this sector. Occupation-tail and square-root
bounds compare its energy and vector with a specified Gaussian ground.
In particular, $0\le E_N-E_{\rm free}\le\epsilon_E$ implies squared
ground overlap at least $1-\epsilon_E/\Delta$.

The quadratic comparison has
\[
 \omega_a(p)^2=\Delta^2\left[1+
 \frac{1-a^2}{3a^2}\sum_i(1-\cos(ap_i))\right]
 \longrightarrow \Delta^2+\frac{\Delta^2}{6}|p|^2.
\]
Complete bounded field instruments retain their defect flags and a
finite orientation reference. Explicit capacity choices bound a
two-instrument, three-evolution record law by $19/20000$, with separate
spacing and volume errors. This is a controlled free Klein--Gordon
description with three internal scalar components. The component labels
do not furnish spatial spin one. Replication, spectral dressing and the
spacing-dependent clock are declared architecture; no empirical mass
or GeV conversion selects their coefficients
(\path{model/interacting_block_propagation},
\path{model/interacting_block_quadrature}, \cite{ophrepo}).

\paragraph{A finite collision distinguishable from free evolution.}
The same paired Hamiltonian has a finite-$N$ channel forbidden by its
free quadratic comparison. Prepare two different internal components
in Fourier packets with disjoint input and output momenta but equal
total momentum. The free evolution preserves each component's momentum;
the capacity square root produces a connected quartic vertex. For four
momenta on one frequency shell,
\[
 \langle k',l'|V_1|k,l\rangle
 =\frac{t_a A(k)}{V}\left[\left(\frac{D}{\omega_k}\right)^2+1\right]>0,
 \quad A(k)=2\sum_i\cos k_i,\quad D=\Delta/a,
\]
when $A(k)>0$, with $H_{a,N}=H_{\rm free}+V_1/N+R_N$.
Global number-weighted bounds control $R_N$, the second Born term,
packet admission and finite preparation before squaring the signal.
The experiment retains its hit, complement and preparation failures.
It uses explicitly prepared nonstationary states; an approximate
interacting ground is not substituted in a rare-event bound.

An independently interval-certified diagnostic on a five-block cycle
evolves the complete $1024$-state reducing carrier, retaining all three
components. With $D/\Delta=2$, $t/\Delta=1/8$ and $\Delta T=1$, its
preparation-inclusive final hit probability is enclosed by
$[0.00296630042443,0.00308409427253]$.
The associated occupation records are quantum marginals, rather than
classical particle paths. At $N=1$, $H_N=P_NH_{\rm free}P_N$;
the admitted packet projection changes the separate momentum supports.
Consequently the diagnostic's nonzero response alone does not certify
a channel forbidden to the free evolution of those same admitted states.
An exact Gaussian-kernel calculation of that matched free evolution,
without an intermediate oscillator cutoff, gives the preparation-inclusive
interval $[0.00218695122151,0.00228825603677]$. The native-minus-free
contrast is therefore at least $0.00067804438766$. Both comparisons retain
the same admitted input, output projector and quantum reference; the
complete native complement has an explicit common-carrier embedding.
The large-$N$ theorem separately bounds that admission error below its
connected interaction signal. Three-dimensional finite members also have
strictly positive controlled collision channels. Their interaction
vanishes as $N\to\infty$ in the established refinement. Neither the
cycle nor those finite records proves a nontrivial interacting
relativistic ultraviolet limit, an asymptotic scattering matrix or a
Standard Model species assignment
(\path{model/interacting_block_scattering}, \cite{ophrepo}).

\paragraph{A family effect and an energy clock.}
For the declared cyclic family register, native repair columns give the
one-tick amplitude $M=\sqrt2\,C(\theta)/3$, where
$C=I/\sqrt6+(e^{i\theta}R+e^{-i\theta}R^{-1})/\sqrt{12}$ and $R^3=I$.
Positive eigenvalues of $C$ have squared values with Koide ratio $2/3$.
The literal two-tick amplitude is instead $\sqrt3 I/9$.
An unconditional collision extension realizes $A=M^\dagger M$ as a
generator: $N$ fresh retained ancillas over time $t$ give marginal-channel
diamond error at most $\min(2,4t^2/N)$ from $e^{-itA}$. The projector
phase gate, inverse tick and schedule are explicit architecture inputs.
The XOR probability $2/9$ does not alone choose a Hamiltonian phase.

The separate scalar-memory construction contains three CAR families and
$5\,2^{18}$ states. On its stationary vacuum-generated component, put
$S=4\sum_f y_f^2$, $D_0=\sqrt{\omega^2+4S}$ and
$D_f=\sqrt{\omega^2+4S-8y_f^2}$. Its lower charged line is
$m_f=4y_f^2/(D_0+D_f)$; a second line, nine dark states and the no-operation
record are retained. For $y_h>y_l>0$, $m_h/m_l>(y_h/y_l)^2$.
Consequently this response cannot remove the $+9.83354$ ppm excess in
the muon/electron comparison of the $\theta=2/9$ squared-amplitude
reading. Three independently occupiable families, a common three-state
register and a neutral classical Higgs background are distinct source
constructions; these finite identities do not transfer their energies to
the spatial interacting source
(\path{model/finite_family_clock}, \cite{ophrepo}).

\paragraph{A defect class and a particle representation.}
The six-class register preserves colour triality and weak centre parity:
$k=3n_2-2n_3\pmod6$. Its $32$ exterior states have class multiplicities
$(4,6,6,4,6,6)$, while the independently computed fifteen-state current
representation gives the internal weights of one Standard Model fermion
family. This comparison requires the full currents. For example, the
centre acts trivially on both summands of
$\mathbf3\otimes\overline{\mathbf3}=\mathbf1\oplus\mathbf8$, whereas the
full quadratic Casimir is $3(I-P_1)$. Two quarks also contain a
$(\overline{\mathbf3},\mathbf1)_{1/3}$ channel with the gauge charges of
$d^c$, but their even fermion parity differs from a single odd fermion.
A class-conserving fusion therefore needs its complete physical reference
and exchange law before receiving a particle interpretation. Auxiliary
exterior degree alone is not the physical fermion grading
(\path{model/particle_catalogue_correspondence}, \cite{ophrepo}).

